\chapter{特征函数与中心极限
定理}\label{ch:b3:clt}

大数定律说平均值收敛；中心极限
定理说的是\emph{它们如何涨落}：误差，
乘以 $\sqrt n$ 后渐近服从\omterm{ex:b3:probability:laws}{高斯分布}——
无论从哪种分布出发。这一普适性是
初等概率论中最深刻的事实，其自然的证明
是傅里叶分析式的：\emph{\omterm{def:b3:clt:cf}{特征函数}}
（即分布的傅里叶变换）把\omterm{def:b3:probability:independence}{独立}和
化为乘积，而\cref{ch:b3:fouriertransform}中的方法
——单射性与高斯不动点——将逐点
收敛的这些乘积转化为分布的收敛
（莱维定理，给出完整证明）。本章最后
以\omterm{def:b3:clt:gaussianvector}{高斯向量}作结，并严格推导统计学中处处使用的置信
区间；周末问题
给出林德伯格的中心极限定理第二种证明，并附明确的
误差阶。

\section{特征函数}

\begin{definition}\label{def:b3:clt:cf}
称实\omterm{def:b3:probability:space}{随机变量} $X$ 的\emph{特征函数}\index{特征函数
}为
\[
\varphi_X(\xi) = \E\bigl[\eu^{\iu\xi X}\bigr]
= \int_\R \eu^{\iu\xi x}\,\dd\P_X(x)
\qquad (\xi \in \R)
\]
（由变换定理从分布算出；若
$f$ 为密度，则 $\varphi_X(\xi) = \hat f(-\xi)$，
沿用\cref{ch:b3:fouriertransform}中的约定）。
\end{definition}

\begin{proposition}\label{prop:b3:clt:cfbasics}
(a) $\varphi_X(0) = 1$，$\abs{\varphi_X} \leq 1$，且
$\varphi_X$ 一致连续；
$\varphi_{aX + b}(\xi) = \eu^{\iu b\xi}\varphi_X(a\xi)$。
(b) 若 $X, Y$ 相互\emph{\omterm{def:b3:probability:independence}{独立}}：
$\varphi_{X+Y} = \varphi_X\,\varphi_Y$。
(c) 若 $\E\abs X^k < \infty$，则 $\varphi_X \in \mathcal
C^k$，且 $\varphi_X^{(j)}(0) = \iu^j\,\E[X^j]$，对 $j \leq
k$；特别地，对中心化的 $X \in L^2$，其方差为
$\sigma^2$：
\[
\varphi_X(\xi) = 1 - \frac{\sigma^2\xi^2}{2} +
o(\xi^2) \qquad (\xi \to 0).
\]
(d) 高斯情形：$X \sim \mathcal N(m, \sigma^2)$ 的\omterm{def:b3:clt:cf}{特征函数}
为 $\varphi_X(\xi) = \eu^{\iu m\xi - \sigma^2\xi^2/2}$。
\end{proposition}

\begin{proof}
(a) 有界性显然；连续性：$\abs{\varphi(\xi + h)
- \varphi(\xi)} \leq \E\abs{\eu^{\iu hX} - 1} \to 0$，当 $h
\to 0$ 时由控制收敛定理得出，且关于 $\xi$ 一致成立。
仿射规则即变量代换。(b) $\eu^{\iu\xi(X+Y)} =
\eu^{\iu\xi X}\eu^{\iu\xi Y}$，且\omterm{def:b3:probability:independence}{独立}变量乘积的\omterm{def:b3:probability:space}{期望}
等于各\omterm{def:b3:probability:space}{期望}之积
（\cref{thm:b3:probability:independence}，分别对实部与
虚部应用）。(c) 对\omterm{def:b3:probability:space}{期望}号下的求导，
由 $\E\abs X^j$ 控制
（\cref{thm:b3:lebesgue:paramdiff}）；
于 $0$ 处的泰勒展开即是对 $\mathcal C^2$ 函数
$\varphi$ 的泰勒–杨公式。(d) 对 $\mathcal N(0,1)$，由高斯变换
（\cref{ex:b3:fouriertransform:gaussian}，取 $a = \frac12$）
得 $\int\eu^{\iu\xi x}\frac{\eu^{-x^2/2}}{\sqrt{2\pi}}\dd
x = \eu^{-\xi^2/2}$；一般情形由仿射规则得到。
\end{proof}

\begin{theorem}[单射性]\label{thm:b3:clt:injectivity}
若 $\varphi_X = \varphi_Y$，则 $X$ 与 $Y$ 有相同的
分布。更确切地说，对 $N \sim \mathcal N(0,1)$（与 $X$ \omterm{def:b3:probability:independence}{独立}）
及 $\varepsilon > 0$，光滑化变量 $X +
\varepsilon N$ 具有如下密度
\[
p_\varepsilon(x) = \frac1{2\pi}\int_\R
\varphi_X(-\xi)\,\eu^{-\varepsilon^2\xi^2/2}\,
\eu^{\iu\xi x}\,\dd\xi ,
\]
仅由 $\varphi_X$ 决定；令 $\varepsilon \to 0$
则恢复 $X$ 的分布。\index{特征函数,
单射性}
\end{theorem}

\begin{proof}
$X + \varepsilon N$ 具有密度 $p_\varepsilon(x) =
\E\bigl[g_\varepsilon(x - X)\bigr]$，其中 $g_\varepsilon$ 为
$\mathcal N(0, \varepsilon^2)$ 的密度：事实上，对博雷尔
$B$，由\omterm{def:b3:probability:independence}{独立}性及托内利定理得 $\P(X + \varepsilon N \in
B) = \int\!\!\int\mathbf 1_B(x + \varepsilon
n)g_1(n)\,\dd n\,\dd\P_X(x) = \int_B\E[g_\varepsilon(t -
X)]\dd t$（先作变量代换，再用一次托内利定理）。
由其变换的傅里叶反演公式写出 $g_\varepsilon$
（\cref{exo:b3:fouriertransform:4}，经伸缩）：
$g_\varepsilon(u) = \frac1{2\pi}\int
\eu^{-\varepsilon^2\xi^2/2}\eu^{\iu\xi u}\dd\xi$，并用富比尼定理
（一切项均由高斯因子控制）：
\[
p_\varepsilon(x) = \frac1{2\pi}\int_\R
\varphi_X(-\xi)\,\eu^{-\varepsilon^2\xi^2/2}\,
\eu^{\iu\xi x}\,\dd\xi ,
\]
仅依赖于 $\varphi_X$ 的泛函。若 $\varphi_X =
\varphi_Y$，则 $X + \varepsilon N$ 与 $Y + \varepsilon N$
对每个 $\varepsilon$ 同分布；对有界连续
$f$，$\E f(X + \varepsilon N) \to \E f(X)$，当 $\varepsilon
\to 0$ 时（由控制收敛定理，$X + \varepsilon N \to X$
在乘积空间上逐点收敛），故 $\E f(X) = \E f(Y)$ 对
所有此类 $f$ 成立——而这决定了分布：对每个 $t$，
将 $\mathbf 1_{\intoc{-\infty}t}$ 夹在有界
连续斜坡函数 $f_k^\pm$（取 $1$ 的范围为
$\intoc{-\infty}{t \mp \frac1k}$，取 $0$ 的范围在 $t \pm
\frac1k$ 之外，其间为仿射）；在 $\E
f_k^-(X) \leq F_X(t) \leq \E f_k^+(X)$ 中取极限，得 $F_X(t) =
F_Y(t)$，对两者都连续的每个 $t$，从而
由右连续性与公共连续点的稠密性，处处成立
（两个 $F$ 的跳点均可数）；
分布函数相等则分布相等
（\cref{exo:b3:measure:3}，依赖于
\cref{thm:b3:measure:uniqueness}）。
\end{proof}

\section{依分布收敛}

\begin{definition}\label{def:b3:clt:cid}
$X_n$ \emph{依分布收敛}（或简称依分布）于 $X$，
记作 $X_n \Rightarrow X$，若\index{依分布
收敛}
\[
\E\bigl[f(X_n)\bigr] \longrightarrow \E\bigl[f(X)\bigr]
\qquad\text{对每个有界连续函数 } f\colon\R\to\R .
\]
等价地（见\cref{exo:b3:clt:4}）：$F_{X_n}(t) \to F_X(t)$
在 $t$ 为 $F_X$ 的每个连续点时成立。$X_n$ 不必
定义在同一\omterm{def:b3:probability:space}{概率空间}上，只有分布才重要。
\end{definition}

\begin{theorem}[赫利选择定理]
\label{thm:b3:clt:helly}
任一分布函数列 $(F_n)$ 都有一
子列，在极限的每个连续点处逐点收敛于
某个非降右连续的 $G \colon
\R \to \intcc01$——可能有 $G(+\infty) - G(-\infty) <
1$（质量可能逃逸至无穷远）。\index{赫利定理}
\end{theorem}

\begin{proof}
由对角线抽取，对每个有理数均有 $F_{n_k}(q) \to \ell(q)$，
（$q$ 处的值取于紧集 $\intcc01$ 中）。定义
$G(t) = \inf\{\ell(q) : q \in \Q, q > t\}$：非降；
右连续（在右侧逐渐收缩的有理
\omterm{def:b3:topology:topology}{邻域}上取下确界）。设 $t$ 为
$G$ 的连续点，对有理数 $q_1 < t < q_2$，
\[
\ell(q_1) \leq \liminf F_{n_k}(t) \leq \limsup F_{n_k}(t)
\leq \ell(q_2),
\]
由各 $F_{n_k}$ 的单调性。由
$G$ 为下确界的定义及 $\ell$ 在有理数上的单调性：
$G(s) \leq \ell(q) \leq G(q)$，只要 $s < q$。
取 $s < q_1 < t$，得 $\ell(q_1) \geq G(s)$，并
$\ell(q_2) \leq G(q_2)$；令 $s \uparrow t$ 及 $q_2
\downarrow t$，由 $G$ 在 $t$ 处的连续性，
$\liminf$ 与 $\limsup$ 都夹向 $G(t)$。
\end{proof}

\begin{lemma}[由特征函数得紧性]
\label{lem:b3:clt:tightness}
对任一\omterm{def:b3:probability:space}{随机变量} $X$ 及 $u > 0$：
\[
\P\Bigl(\abs X \geq \frac2u\Bigr) \;\leq\;
\frac1u\int_{-u}^{u}\bigl(1 -
\operatorname{Re}\varphi_X(\xi)\bigr)\,\dd\xi .
\]
\end{lemma}

\begin{proof}
由托内利--富比尼定理（被积函数有界，积分区域在
$\xi$ 上有界）：
\[
\frac1u\int_{-u}^u\bigl(1 -
\operatorname{Re}\varphi_X(\xi)\bigr)\dd\xi
= \E\Bigl[\frac1u\int_{-u}^u(1 - \cos(\xi X))\,\dd\xi\Bigr]
= 2\,\E\Bigl[1 - \frac{\sin(uX)}{uX}\Bigr]
\]
（将方括号理解为 $0$，即它在 $X = 0$ 处的极限）。
被积函数非负（$\abs{\sin t} \leq \abs t$），且
当 $\abs{uX} \geq 2$ 时：$1 - \frac{\sin(uX)}{uX} \geq 1 -
\frac1{\abs{uX}} \geq \frac12$。仅保留\omterm{def:b3:probability:space}{期望}中的
事件 $\{\abs{uX} \geq 2\}$，因此
至少留下 $2 \cdot \frac12\,\P(\abs X \geq \frac2u)$，
即为所证。
\end{proof}

\begin{theorem}[莱维连续性定理]
\label{thm:b3:clt:levy}
设 $(X_n)$ 为\omterm{def:b3:probability:space}{随机变量}序列，其\omterm{def:b3:clt:cf}{特征函数}
逐点收敛：$\varphi_{X_n}(\xi) \to
\varphi(\xi)$，对每个 $\xi$ 成立，其中 $\varphi =
\varphi_X$ 为某个\omterm{def:b3:probability:space}{随机变量}
$X$ 的\omterm{def:b3:clt:cf}{特征函数}。则 $X_n \Rightarrow X$。\index{莱维
连续性定理}
\end{theorem}

\begin{proof}
\emph{紧性。}取 $\varepsilon > 0$。由于 $\varphi$ 在
$0$ 处连续，且 $\varphi(0) = 1$，取 $u > 0$ 使
$\frac1u\int_{-u}^u(1 - \operatorname{Re}\varphi) <
\varepsilon$；由控制收敛定理（被积函数以
$2$ 为界（区间 $[-u,u]$ 固定），对
$\varphi_{X_n}$ 的同一积分 $< 2\varepsilon$（$n$ 充分大时）：
由\cref{lem:b3:clt:tightness}得 $\P(\abs{X_n} \geq \frac2u)
\leq 2\varepsilon$（$n$ 充分大时），并通过增大常数
处理其余有限多个：各分布是\emph{紧的}
——没有质量逃逸。

\emph{子列。}设 $(F_{n_k})$ 为任一子列，
由\cref{thm:b3:clt:helly}，可抽取 $F_{n_{k_j}} \to G$
（在连续点处）。紧性迫使 $G(-\infty) = 0$，
$G(+\infty) = 1$（$G(\frac2u) - G(-\frac2u) \geq 1 -
2\varepsilon$，在连续点处成立）：$G$ 是真正的分布函数，
即某\omterm{def:b3:probability:space}{随机变量} $Y$ 的分布函数。于是
$X_{n_{k_j}} \Rightarrow Y$（见\cref{exo:b3:clt:4}，
由 $F$ 的收敛即\omterm{def:b3:clt:cid}{依分布收敛}），故
$\varphi_{X_{n_{k_j}}} \to \varphi_Y$ \emph{逐点}成立（$x
\mapsto \eu^{\iu\xi x}$ 是有界连续函数，实部与
虚部分别处理）；与假设比较得
$\varphi_Y = \varphi = \varphi_X$，并由单射性
（见\cref{thm:b3:clt:injectivity}）得 $Y \sim X$，即 $G =
F_X$。

\emph{结论。}$(F_n)$ 的每个子列都有一子子列，
收敛于\emph{同一}个 $F_X$（在其
连续点上）；因此 $F_n(t) \to F_X(t)$ 对每个
连续点 $t$ 成立（实数列若其每个子列都有
具有相同极限的子子列，则该数列收敛）：
故 $X_n \Rightarrow X$。
\end{proof}

\section{中心极限定理}

\begin{theorem}[中心极限定理]\label{thm:b3:clt:clt}
设 $(X_n)$ 为\omterm{def:b3:probability:independence}{独立}同分布，$\E X_1 = m$ 且 $\V(X_1) =
\sigma^2 \in \intoo0\infty$，则\index{中心
极限定理}
\[
\frac{S_n - nm}{\sigma\sqrt n} \;\Longrightarrow\; \mathcal
N(0, 1) :
\qquad
\P\Bigl(a \leq \frac{S_n - nm}{\sigma\sqrt n} \leq
b\Bigr) \longrightarrow
\frac{1}{\sqrt{2\pi}}\int_a^b\eu^{-x^2/2}\,\dd x
\]
对所有 $a < b$。
\end{theorem}

\begin{proof}
中心化并标准化：$Z_i = \frac{X_i - m}{\sigma}$
（\omterm{def:b3:probability:independence}{独立}同分布，均值为 $0$，方差为 $1$），并记 $T_n =
\frac1{\sqrt n}\sum_{i\leq n}Z_i$。由\omterm{def:b3:probability:independence}{独立}性及
仿射性质（\cref{prop:b3:clt:cfbasics}）：
\[
\varphi_{T_n}(\xi) =
\varphi_{Z}\Bigl(\frac{\xi}{\sqrt n}\Bigr)^{n},
\qquad
\varphi_Z(\eta) = 1 - \frac{\eta^2}2 + \eta^2\rho(\eta),\quad
\rho(\eta)\to0 .
\]
取定 $\xi$，令 $a_n = \varphi_Z(\xi/\sqrt n)$，$b_n = 1 -
\frac{\xi^2}{2n}$，二者的模均 $\leq 1$（只要 $n$ 充分大）；
$\abs{b_n} \leq 1$ 当 $\xi^2 \leq 4n$ 时成立，$\abs{a_n} \leq 1$
总成立。由初等不等式 $\abs{a^n - b^n} \leq
n\abs{a - b}$，对 $\abs a, \abs b \leq 1$ 成立（裂项相消：
$a^n - b^n = \sum a^k(a - b)b^{n-1-k}$），可得
\[
\Bigl|\varphi_{T_n}(\xi) - \Bigl(1 -
\frac{\xi^2}{2n}\Bigr)^{n}\Bigr|
\leq n\,\Bigl|\varphi_Z\Bigl(\frac\xi{\sqrt n}\Bigr) - 1 +
\frac{\xi^2}{2n}\Bigr|
= \xi^2\,\Bigl|\rho\Bigl(\frac{\xi}{\sqrt n}\Bigr)\Bigr|
\longrightarrow 0,
\]
而 $\bigl(1 - \frac{\xi^2}{2n}\bigr)^n \to
\eu^{-\xi^2/2}$（取实对数）。于是 $\varphi_{T_n}(\xi) \to
\eu^{-\xi^2/2} = \varphi_{\mathcal N(0,1)}(\xi)$
（\cref{prop:b3:clt:cfbasics}(d)）对每个 $\xi$ 成立。由列维
连续性定理（\cref{thm:b3:clt:levy}）推出 $T_n \Rightarrow \mathcal
N(0,1)$。区间概率随之成立，因为 $F_{\mathcal
N}$ 处处连续。
\end{proof}

\begin{example}[置信区间，诚实推导]
\label{ex:b3:clt:confidence}
对 $n$ 位\omterm{def:b3:probability:independence}{独立}选民作民意调查；$\hat p_n = S_n/n$ 用于估计
真实比例 $p$，其中 $\sigma^2 = p(1-p) \leq \frac14$。由中心极限定理
可得，对充分大的 $n$，
\[
\P\Bigl(\abs{\hat p_n - p} \leq
\frac{z}{2\sqrt n}\Bigr)
\;\geq\; \P\Bigl(\Bigl|\frac{S_n - np}{\sigma\sqrt n}\Bigr|
\leq z\Bigr)
\longrightarrow \Phi(z) - \Phi(-z),
\]
其中 $\Phi$ 为标准\omterm{ex:b3:probability:laws}{高斯分布}的分布函数。
取 $z = 1.96$，此时渐近置信水平为 $95\%$，而
误差限 $\frac{1.96}{2\sqrt n} \leq 3\%$ 要求 $n \geq
\bigl(\frac{1.96}{0.06}\bigr)^2 \approx 1068$——这正是
人们在每条“$\pm3$ 个百分点，$95\%$”报道背后读到的数字；与
切比雪夫给出的 $5556$（\cref{exo:b3:probability:7}）相比。
$\sqrt n$ 因子具有普遍性：要使误差减半，
须将样本量增至四倍——这与决定\omterm{ex:b3:probability:sllnapps}{蒙特卡罗}方法代价的规律相同
（\cref{exo:b3:clt:7}）。
\end{example}

\section{高斯向量}

\begin{definition}\label{def:b3:clt:gaussianvector}
称随机向量 $X = (X_1, \dots, X_d)$ 为
\emph{高斯向量}\index{高斯向量}，若其每个线性组合
$\langle t, X\rangle = \sum t_iX_i$ 均为
（可能退化的）实高斯变量。其分布
由均值向量 $m = (\E X_i)$ 及
\emph{协方差矩阵} $\Sigma = \bigl(\operatorname{Cov}
(X_i, X_j)\bigr)$：事实上，
向量的\omterm{def:b3:clt:cf}{特征函数} $\varphi_X(t) = \E\eu^{\iu\langle t, X\rangle}$ 恰为
在 $1$ 处取值的 $\langle t, X\rangle$ 的\omterm{def:b3:clt:cf}{特征函数}：
\[
\varphi_X(t) = \exp\Bigl(\iu\langle t, m\rangle -
\tfrac12\,t^{\mathsf T}\Sigma\,t\Bigr),
\]
且 $d$ 维\omterm{def:b3:clt:cf}{特征函数}是单射
（证明与\cref{thm:b3:clt:injectivity}中的光滑化方法相同，
逐坐标取高斯变量）。
\end{definition}

\begin{theorem}\label{thm:b3:clt:gaussianvector}
设 $X$ 为\omterm{def:b3:clt:gaussianvector}{高斯向量}。
\begin{enumerate}
  \item 任一仿射像 $AX + b$ 均为\omterm{def:b3:clt:gaussianvector}{高斯向量}。
  \item 分量 $X_i$ 相互\emph{\omterm{def:b3:probability:independence}{独立}}当且仅当
        $\Sigma$ 为对角矩阵：对联合高斯变量，
        不相关 $=$ \omterm{def:b3:probability:independence}{独立}。
  \item 若 $\Sigma$ 可逆，则 $X$ 具有密度
        $\frac{1}{(2\pi)^{d/2}\sqrt{\det\Sigma}}
        \exp\bigl(-\frac12(x - m)^{\mathsf T}\Sigma^{-1}(x -
        m)\bigr)$。
\end{enumerate}
\end{theorem}

\begin{proof}
(1) $AX + b$ 各分量的线性组合，
均为 $X$ 的线性组合的仿射函数，而高斯变量的仿射像仍是高斯的，
故此向量为\omterm{def:b3:clt:gaussianvector}{高斯向量}。
(2) 若 $\Sigma$ 为对角矩阵，则\omterm{def:b3:clt:cf}{特征函数}
可分解：$\varphi_X(t) = \prod_i\exp(\iu t_im_i -
\frac12\Sigma_{ii}t_i^2) = \prod\varphi_{X_i}(t_i)$，此即
乘积分布的\omterm{def:b3:clt:cf}{特征函数}
（由\cref{thm:b3:probability:independence}，
借 $d$ 维单射性读出）：各分量相互\omterm{def:b3:probability:independence}{独立}。
其逆为：\omterm{def:b3:probability:independence}{独立}的
$L^2$ 变量之协方差为零。
(3) 对 $\Sigma = P D P^{\mathsf T}$ 作对角化（$P$ 为正交矩阵，
$D > 0$ 为对角矩阵——见\cref{exo:b3:submanifolds:8}）；
向量 $Y = P^{\mathsf T}(X - m)$ 为\omterm{def:b3:clt:gaussianvector}{高斯向量}，其协方差为
$D$：由 (2)，其分量相互\omterm{def:b3:probability:independence}{独立}，
各分量服从 $\mathcal N(0, d_i)$，故 $Y$ 的密度为乘积形式；
再经保体积变换 $x = m + PY$ 作推前
（由\cref{thm:b3:product:linearchange}，$\abs{\det P} = 1$），并
将指数改写为不依赖坐标的形式。
\end{proof}

\begin{theorem}[多维中心极限定理]
\label{thm:b3:clt:multiclt}
设 $(X_n)$ 为\omterm{def:b3:probability:independence}{独立}同分布的平方\omterm{def:b3:lebesgue:l1}{可积}随机
\emph{向量}，取值于 $\R^d$，均值为 $m$，协方差矩阵为
$\Sigma$。则 $\frac{S_n - nm}{\sqrt n}$ \omterm{def:b3:clt:cid}{依分布收敛}于
\omterm{def:b3:clt:gaussianvector}{高斯向量} $\mathcal N(0,
\Sigma)$。\index{中心极限定理!多维}
\admitted
\end{theorem}

\begin{remark}\label{rem:b3:clt:multiclt}
几乎全部材料已在手边。对每个
方向 $t \in \R^d$，实\omterm{def:b3:probability:space}{随机变量} $\langle t,
\frac{S_n - nm}{\sqrt n}\rangle$ 是
方差为 $t^{\mathsf T}\Sigma t$ 的\omterm{def:b3:probability:independence}{独立}同分布实\omterm{def:b3:probability:space}{随机变量}的规范化和，
故由\cref{thm:b3:clt:clt}的计算得
$d$ 维\omterm{def:b3:clt:cf}{特征函数}逐点收敛，
趋于 $\eu^{-t^{\mathsf T}\Sigma t/2}$，即
$\mathcal N(0, \Sigma)$ 的\omterm{def:b3:clt:cf}{特征函数}
（见\cref{def:b3:clt:gaussianvector}）。尚未重新证明的是
\emph{在 $\R^d$ 中}的莱维连续性定理：
赫利选择定理与紧性估计可逐坐标推广，
这是常规的（逐坐标进行），而此处的克拉默–沃尔德化归
在任何研究生概率课程中都会认真完成；
此外无需超出本章方法的
内容。
\end{remark}

\begin{method}\label{met:b3:clt:method}
欲识别极限分布：计算\omterm{def:b3:clt:cf}{特征函数}，
取逐点极限，辨认其形式（高斯
$\eu^{-\sigma^2\xi^2/2}$，泊松
$\eu^{\lambda(\eu^{\iu\xi}-1)}$，指数分布 $\frac{\lambda}
{\lambda - \iu\xi}$，\dots），并援引莱维定理。三步
程式（\omterm{def:b3:probability:independence}{独立}性 $\to$ 乘积；$0$ 处泰勒展开 $\to$
指数极限；莱维 $\to$ \omterm{def:b3:clt:cid}{依分布收敛}）证明
中心极限定理、稀有事件的泊松分布
（见\cref{exo:b3:clt:5}），以及本课程中每一个
经典极限定理。对几乎必然成立的陈述，请回到
\cref{ch:b3:probability}的工具：两章回答
关于同一 $S_n$ 的不同问题。
\end{method}

\section{练习}

\begin{exercise}[$\star$]\label{exo:b3:clt:1}
求下列分布的\omterm{def:b3:clt:cf}{特征函数}：
$\intcc{-1}1$ 上的均匀分布；指数分布 $\mathcal E(\lambda)$；
泊松分布 $\mathcal P(\lambda)$；二项分布 $\mathcal B(n, p)$。
由\cref{thm:b3:clt:injectivity}可知，\omterm{def:b3:probability:independence}{独立}的
泊松变量（参数 $\lambda, \mu$）之和服从泊松分布 $(\lambda +
\mu)$。
\end{exercise}

\begin{exercise}[$\star\star$]\label{exo:b3:clt:2}
(a) 证明：$\varphi_X$ 为实值，当且仅当 $X$
与 $-X$ 同分布（称为\emph{对称}变量）。
(b) 设 $\abs{\varphi_X(\xi_0)} = 1$，其中 $\xi_0 \neq
0$，证明 $X$ 几乎必然支撑于等差数列
$a + \frac{2\pi}{\xi_0}\Z$ 上（\emph{记
$\varphi_X(\xi_0) = \eu^{\iu\theta}$，并计算 $\E[1 -
\cos(\xi_0X - \theta)]$）}。由此可得，若 $X$ 有密度，
则 $\abs{\varphi_X(\xi)} < 1$ 对一切 $\xi \neq 0$ 成立。
\end{exercise}

\begin{exercise}[$\star\star$]\label{exo:b3:clt:3}
设 $X \sim \mathcal N(m_1, \sigma_1^2)$，$Y \sim
\mathcal N(m_2, \sigma_2^2)$ 相互\omterm{def:b3:probability:independence}{独立}。证明 $X + Y
\sim \mathcal N(m_1 + m_2, \sigma_1^2 + \sigma_2^2)$，并
更一般地，证明高斯族在\omterm{def:b3:probability:independence}{独立}和与仿射映射下
保持稳定。对比：两个\emph{相依}高斯变量之和
是否总是高斯的？
（见\cref{exo:b3:clt:9}。）
\end{exercise}

\begin{exercise}[$\star\star$]\label{exo:b3:clt:4}
(a) 证明\cref{def:b3:clt:cid}中的等价性：若
$\E f(X_n) \to \E f(X)$ 对一切有界连续函数 $f$ 成立，
则 $F_{X_n}(t) \to F_X(t)$ 在连续点处成立，
\emph{（用两个连续阶梯状斜坡函数夹住
$\mathbf 1_{\intoc{-\infty}t}$）}；反之
\emph{（用斜坡函数之和逼近有界连续 $f$，
或在连续点的细网格上作条件化）}
——逆命题可先对一致连续的 $f$ 处理，
再推广到一般情形。
(b) 证明：若 $X_n \Rightarrow c$（常数），则 $X_n
\to c$ 依概率收敛。
\end{exercise}

\begin{exercise}[$\star\star$]\label{exo:b3:clt:5}
（稀有事件定律）设 $X_n \sim \mathcal B(n, p_n)$，其中
$np_n \to \lambda > 0$。利用\omterm{def:b3:clt:cf}{特征函数}
及\cref{thm:b3:clt:levy}，证明 $X_n \Rightarrow \mathcal
P(\lambda)$。数值检验：比较 $\P(X = 0)$
在 $\mathcal B(100, 0.02)$ 与 $\mathcal P(2)$ 下的取值。
\end{exercise}

\begin{exercise}[$\star\star$]\label{exo:b3:clt:6}
(a) 投掷一颗均匀骰子 $n = 1000$ 次，近似计算
总点数超过 $3600$ 的概率（均值为 $3500$，
每次投掷的方差为 $\frac{35}{12}$）。
(b) 设 $S \sim \mathcal B(100, \frac12)$，用
带连续性修正的中心极限定理近似计算 $\P(45 \leq S \leq 55)$，
其中修正量取 $\pm\frac12$，
并评述该修正的作用。
\end{exercise}

\begin{exercise}[$\star\star$]\label{exo:b3:clt:7}
（\omterm{ex:b3:probability:sllnapps}{蒙特卡罗}误差）
在\cref{pb:b3:probability:1}第 11 问的设定下，取 $g \in
L^2(\intcc01^d)$，记 $\sigma^2 = \V(g(U_1))$，且 $I = \int
g$。证明
\[
\sqrt n\,\Bigl(\frac1n\sum_{k\leq n}g(U_k) - I\Bigr)
\Longrightarrow \mathcal N(0, \sigma^2),
\]
并由此推出渐近 $95\%$ 误差界 $\pm
1.96\,\sigma/\sqrt n$——与维数 $d$ 无关。
与确定性中点法则相比较，该法则在
$d$ 维下误差为 $\sim n^{-2/d}$（对 $\mathcal C^2$ 被积函数）：
从哪一维数起，随机抽样更优？
\end{exercise}

\begin{exercise}[$\star\star\star$]\label{exo:b3:clt:8}
（斯卢茨基）设 $X_n \Rightarrow X$，且 $Y_n \to c$ 依
概率收敛（$c$ 为常数）。证明 $X_n + Y_n \Rightarrow X +
c$，且 $Y_nX_n \Rightarrow cX$。\emph{（利用
\omterm{def:b3:clt:cf}{特征函数}和不等式 $\abs{\E\eu^{\iu\xi
(X_n+Y_n)} - \eu^{\iu\xi c}\E\eu^{\iu\xi X_n}} \leq
\E\abs{\eu^{\iu\xi(Y_n - c)} - 1}$，按 $\abs{Y_n - c}
\leq \delta$ 分两种情形讨论。）} 应用：
在\cref{ex:b3:clt:confidence}中，论证将未知量
$\sigma = \sqrt{p(1-p)}$ 换为 $\sqrt{\hat p_n(1 - \hat
p_n)}$ 的合理性。
\end{exercise}

\begin{exercise}[$\star\star\star$]\label{exo:b3:clt:9}
设 $X \sim \mathcal N(0,1)$，$\varepsilon$ 与之\omterm{def:b3:probability:independence}{独立}，
且 $\P(\varepsilon = \pm1) = \frac12$；令 $Y =
\varepsilon X$。
(a) 证明 $Y \sim \mathcal N(0,1)$ 且
$\operatorname{Cov}(X, Y) = 0$。
(b) 证明 $X$ 与 $Y$ \emph{不}\omterm{def:b3:probability:independence}{独立}，并说明
$(X, Y)$ 不是\omterm{def:b3:clt:gaussianvector}{高斯向量} \emph{（计算 $\P(X + Y =
0)$）}。
(c) 寓意：\cref{thm:b3:clt:gaussianvector}(2) 要求
联合高斯性——仅凭“不相关的高斯变量”本身，
不能说明任何问题。
\end{exercise}

\begin{exercise}[$\star\star$]\label{exo:b3:clt:10}
柯西分布的密度为 $\frac1{\pi(1 + x^2)}$。
(a) 证明其\omterm{def:b3:clt:cf}{特征函数}为 $\eu^{-\abs\xi}$
（利用\cref{exo:b3:fouriertransform:1} 及反演公式）。
(b) 证明：若 $X_1, \dots, X_n$ 为\omterm{def:b3:probability:independence}{独立}同分布的柯西变量，则
$\frac{S_n}n$ 仍服从柯西分布——且为\emph{同一}分布：
平均值永不集中。
(c) 试与大数定律和中心极限定理相协调：
哪些假设不成立？（计算 $\E\abs{X_1}$。）
\end{exercise}

\begin{exercise}[$\star\star$]\label{exo:b3:clt:11}
（稳定律的雏形）设 $(X_n)$ 为\omterm{def:b3:probability:independence}{独立}同分布的
标准柯西变量（\cref{exo:b3:clt:10}）。
(a) 证明：对任意 $a, b > 0$，$aX_1 + bX_2$ 的分布
与 $(a + b)X_1$ 同分布：柯西族是\emph{严格
稳定}的，指数为 $1$。
(b) 证明高斯族是严格稳定的，指数为
$2$：$aX_1 + bX_2 \sim \sqrt{a^2 + b^2}\,X_1$，对
$X_i$ \omterm{def:b3:probability:independence}{独立}同分布，服从 $\mathcal N(0,1)$。
(c) 利用形如
$\eu^{-c\abs\xi^\alpha}$ 的\omterm{def:b3:clt:cf}{特征函数}，解释为何指数为 $\alpha$ 的稳定性要求
和的规范化为 $n^{1/\alpha}$，并说明
这对中心极限定理的吸引域意味着什么：哪些
\omterm{def:b3:probability:independence}{独立}同分布的和经仿射规范化后，可收敛于
柯西分布而非\omterm{ex:b3:probability:laws}{高斯分布}？
\end{exercise}

\begin{exercise}[$\star\star$]\label{exo:b3:clt:12}
（经验分布函数）设 $(X_n)$ 为\omterm{def:b3:probability:independence}{独立}同分布的
\omterm{def:b3:probability:space}{随机变量}，分布函数为 $F$，记 $F_n(t) =
\frac1n\#\{k \leq n : X_k \leq t\}$。
(a) 固定 $t$。证明 $n F_n(t) \sim \mathcal B(n, F(t))$，
$F_n(t) \to F(t)$ 几乎必然成立（\cref{thm:b3:probability:slln}），
并证明
\[
\sqrt n\,\bigl(F_n(t) - F(t)\bigr) \Longrightarrow
\mathcal N\bigl(0,\ F(t)(1 - F(t))\bigr) .
\]
(b) 渐近方差在哪个 $t$ 处最大？
解释：中位数恰是经验分布
最难确定的位置。
(c) 设 $F$ 连续，证明
$\sup_t\abs{F_n(t) - F(t)}$ 的分布与 $F$ 无关
\emph{（化为均匀变量，借助
\cref{exo:b3:probability:1}）}——这一
与分布无关的奇迹，正是柯尔莫哥洛夫--斯米尔诺夫检验的基础；此处
不要求计算该分布。
\end{exercise}

\section{问题：林德伯格关于中心极限定理的证明，附
带收敛速度}

\begin{problem}[{周末问题——替换法}]
\label{pb:b3:clt:1}
林德伯格（1922 年）借助一个极为朴素的想法
证明了中心极限定理：\emph{逐个把求和项换成高斯变量}，
并用泰勒展开控制每一次替换。该方法无需傅里叶分析，
能给出显式的误差阶，
如今已成为概率论中普适性证明的有力工具。
设 $(X_i)$ 为\omterm{def:b3:probability:independence}{独立}同分布、均值为零的\omterm{def:b3:probability:space}{随机变量}，
$\V(X_1) = 1$，且 $\beta = \E\abs{X_1}^3 < \infty$；设
$(N_i)$ 为\omterm{def:b3:probability:independence}{独立}同分布的 $\mathcal N(0,1)$ 变量，且与
$X_i$ 相互\omterm{def:b3:probability:independence}{独立}（存在性见\cref{thm:b3:probability:existence}）。记
\[
T_n = \frac{X_1 + \dots + X_n}{\sqrt n},
\qquad
G_n = \frac{N_1 + \dots + N_n}{\sqrt n} \sim \mathcal N(0,1).
\]

\medskip
\noindent\textbf{第一部分——交换恒等式。}取定 $f
\in \mathcal C^3_b(\R)$（即三阶导数均有界且连续；
$M_3 = \sup\abs{f'''}$）。对 $0 \leq i \leq n$，
定义混合和
\[
H_i = \frac{X_1 + \dots + X_i + N_{i+1} + \dots +
N_n}{\sqrt n},
\]
于是 $H_n = T_n$，$H_0 = G_n$。
\begin{enumerate}
  \item 写出 $H_i = W_i + \frac{X_i}{\sqrt n}$ 与 $H_{i-1}
        = W_i + \frac{N_i}{\sqrt n}$，其中 $W_i =
        \frac{1}{\sqrt n}\bigl(\sum_{j<i}X_j +
        \sum_{j>i}N_j\bigr)$，注意到 $W_i$
        与数对 $(X_i, N_i)$ \omterm{def:b3:probability:independence}{独立}，请说明理由。
  \item 利用带积分余项或拉格朗日余项的泰勒公式，对任意
        实数 $w, h$，有：
        \[
        \Bigl|f(w + h) - f(w) - f'(w)h -
        \tfrac12f''(w)h^2\Bigr| \leq
        \frac{M_3\,\abs h^3}{6} .
        \]
  \item 两次应用第 2 问（分别取 $h = \frac{X_i}{\sqrt n}$ 与
        $h = \frac{N_i}{\sqrt n}$，并令 $w = W_i$），两边取\omterm{def:b3:probability:space}{期望}，
        再利用\omterm{def:b3:probability:independence}{独立}性，以及
        $X_i$ 与 $N_i$ 前两阶矩相匹配，证明
        \[
        \bigl|\E f(H_i) - \E f(H_{i-1})\bigr|
        \leq \frac{M_3}{6}\cdot
        \frac{\beta + \gamma}{n^{3/2}},
        \qquad \gamma = \E\abs{N_1}^3 =
        \frac{2\sqrt2}{\sqrt\pi} .
        \]
  \item 对 $i$ 裂项求和，并得到\emph{林德伯格
        界}：
        \[
        \bigl|\E f(T_n) - \E f(G_n)\bigr| \leq
        \frac{M_3\,(\beta + \gamma)}{6\,\sqrt n} .
        \]
\end{enumerate}

\medskip
\noindent\textbf{第二部分——从光滑函数 $f$ 到中心极限定理。}
\begin{enumerate}[resume]
  \item 证明 $\E f(T_n) \to \E f(N)$，对每个 $f \in
        \mathcal C_b^3$ 成立，并把结论推广到所有有界
        连续函数 $f$：给定此类 $f$ 与 $\varepsilon$，
        构造 $f_\varepsilon \in \mathcal C^3_b$，使得
        $\norm{f - f_\varepsilon}_\infty \leq \varepsilon$
        在一大区间上满足该估计——例如将 $f$ 与一个
        $\mathcal C^\infty$ 光滑\omterm{def:b3:lp:mollifier}{凸起函数}作卷积
        （见\cref{thm:b3:lp:regularization}）——并借助紧性处理
        尾部（$\V(T_n) = 1$ 与切比雪夫不等式）。
        由此得 $T_n \Rightarrow \mathcal N(0, 1)$：这就是
        重新证明的中心极限定理。
  \item 证明中何处用到了 $X_i$ 的
        \emph{同分布}性质？请说明实际上几乎没有用到它：
        叙述并证明：对\omterm{def:b3:probability:independence}{独立}、均值为零、非同分布的
        $X_i$ 的版本，
        设 $\sum_i\V(X_i) = s_n^2$，并假设三阶矩存在，
        得到误差 $\frac{M_3}{6s_n^3}\sum_i\bigl(\E\abs{X_i}^3
        + \V(X_i)^{3/2}\gamma\bigr)$——林德伯格的真正定理，
        即其李雅普诺夫形式。
\end{enumerate}

\medskip
\noindent\textbf{第三部分——定量收获。}
\begin{enumerate}[resume]
  \item （分布函数）设 $t \in \R$，
        以 $\mathbf 1_{\intoc{-\infty}t}$ 为对象，从上、下两侧
        用 $\mathcal C^3_b$ 斜坡函数逼近，斜坡宽度为 $\delta$
        （构造它们，使 $M_3 = O(\delta^{-3})$）。
        结合第一部分，导出如下双项界：
        \[
        \sup_{t\in\R}\,\bigl|\P(T_n \leq t) -
        \Phi(t)\bigr| \;\leq\;
        \frac{C_1(\beta + \gamma)}{\delta^3\sqrt n} +
        C_2\,\delta
        \qquad (\text{对每个 } \delta > 0),
        \]
        常数可显式给出（$C_2\delta$ 一项用到了
        $\Phi$ 的密度有界，界为
        $\frac1{\sqrt{2\pi}}$），并取 $\delta \sim
        n^{-1/8}$，从而得到一致的
        阶为 $n^{-1/8}$ 的速率。
        （最优阶 $n^{-1/2}$ 即贝里–埃森结果——
        需要更精细的工具；要点在于由初等交换法得到一个\emph{显式}
        速率。）
  \item （定量的棣莫弗–拉普拉斯）特化到
        $X_i = 2B_i - 1$（公平硬币的符号）：将
        结论与局部估计比较（见
        \cref{pb:b3:product:1} 第 7 问）——两种方法各自给出了
        另一种方法所没有的什么？
  \item （普适性）用一段话解释，为何替换法所证明的内容超过
        中心极限定理：任何
        形如 $\E f(\text{和})$ 的统计量，其中 $f$ 光滑，
        在 $n^{-1/2}$ 阶上对求和项的
        \emph{整个分布}（超出前两阶矩的部分）不敏感——
        “不变性原理”，它
        构成了现代普适性结果（如随机矩阵、随机多项式）的基础，
        而中心极限定理是其中
        最早的实例。
\end{enumerate}

\medskip
\noindent\textbf{第四部分——推进的光滑化：更好的
速率。}从 $n^{-1/2}$（光滑 $f$）到
$n^{-1/8}$（分布函数）的损失，源于在上确界范数下
估计 $f'''$。混合量可以弥补其中一部分：它们
含有高斯求和项，而高斯变量具有\emph{光滑化}作用。
\begin{enumerate}[resume]
  \item （隐藏的高斯项）对于 $1 \leq i \leq n - 1$，$h =
        \frac{X_i}{\sqrt n}$ 或 $\frac{N_i}{\sqrt n}$，且
        $\theta \in \intcc01$，写成 $W_i + \theta h = A +
        Z$，其中 $Z = \frac{N_{i+1} + \dots + N_n}{\sqrt
        n}$。证明 $Z \sim \mathcal N\bigl(0,
        \frac{n-i}n\bigr)$ 与数对 $(A,
        h)$ \omterm{def:b3:probability:independence}{独立}，并由此推出，对每个连续函数 $g \in
        L^1(\R)$，
        \[
        \E\bigl[\abs h^3\,\abs{g(W_i + \theta h)}\bigr]
        \;\leq\; \sqrt{\frac{n}{2\pi(n - i)}}\;
        \norm{g}_{L^1}\;\E\abs h^3 .
        \]
  \item 结合第 10 题与积分形式的
        泰勒余项，
        \[
        f(w + h) = f(w) + f'(w)h + \tfrac12f''(w)h^2 +
        \int_0^1\frac{(1 - \theta)^2}2\,f'''(w + \theta
        h)\,h^3\,\dd\theta,
        \]
        用于重做第 3--4 题：设 $f \in \mathcal C^3_b$，
        且另有 $f''' \in L^1(\R)$，
        \[
        \bigl|\E f(T_n) - \E f(G_n)\bigr| \leq
        \frac{\beta + \gamma}{3\sqrt{2\pi}}\cdot
        \frac{\norm{f'''}_{L^1}}{\sqrt n} +
        \frac{M_3(\beta + \gamma)}{6\,n^{3/2}}
        \]
        \emph{（第 10 题处理交换 $i \leq n - 1$
——用 $\sum_{m=1}^{n-1}m^{-1/2} \leq 2\sqrt n$
——而第 3 题的粗略界处理最后一个
        换位）}。验证第 7 题中的斜坡函数满足
        $\norm{\psi_\delta'''}_{L^1} = K_1\delta^{-2}$
        而 $M_3 = K\delta^{-3}$，将其代入，并
        优化 $\delta$：一致的
        分布函数速率改进为
        $O(n^{-1/6})$。
  \item （匹配再多一阶矩）另设 $\E
        X_1^3 = 0$ 且 $\beta_4 = \E X_1^4 < \infty$。
        计算 $\E N_1^3$ 与 $\E N_1^4$，展开到四阶，
        并沿同样的思路证明，分布函数的
        速率变为 $O(n^{-1/4})$
        \emph{（此时 $\norm{\psi_\delta^{(4)}}_{L^1} =
        K_2\delta^{-3}$ 且 $M_4 = K'\delta^{-4}$；取
        $\delta = n^{-1/4}$）}。
  \item （障碍）假设前 $k$ 阶矩
        $X_1$ 的矩与高斯情形一致（$k = 2$ 总成立；
        $k = 3$ 当且仅当 $\E X_1^3 = 0$ 时成立；$k \geq 4$
        几乎从不成立（因为 $\E N_1^4 = 3$）。验证
        第 10--12 题的方案给出分布函数
        的速率
        $n^{-(k-1)/(2k+2)}$，这通过平衡
        $\delta^{-k}n^{-(k-1)/2}$ 与 $\delta$ 得到，并
        观察到，该指数趋于
        贝里--埃森值 $\frac12$，仅当 $k \to
        \infty$ 时如此。用几句话说明，为何
        换位法会饱和：每次换位都按绝对值计入，
        而傅里叶路线
        （埃森光滑化不等式）利用
        \omterm{def:b3:clt:cf}{特征函数}之差的振荡，
        并达到 $C\beta n^{-1/2}$，仅需三阶矩
        即可。
\end{enumerate}

\medskip
\noindent\textbf{第五部分——二维情形：
多维中心极限定理，用换位法。}现设 $X_i$ 为
\omterm{def:b3:probability:independence}{独立}同分布的中心化随机\emph{向量}，取值于 $\R^2$，协方差矩阵为
$\Sigma$，且 $\beta' = \E\norm{X_1}^3 <
\infty$（欧几里得范数）。
\begin{enumerate}[resume]
  \item （\omterm{def:b3:clt:gaussianvector}{高斯向量}，精确地）将 $\Sigma =
        PDP^{\mathsf T}$ 对角化（见\cref{exo:b3:submanifolds:8}），并
        令 $C = P\sqrt DP^{\mathsf T}$。设 $Z = (Z^1,
        Z^2)$ 为一对\omterm{def:b3:probability:independence}{独立}的标准高斯变量
        （见\cref{thm:b3:probability:existence}），证明 $N
        = CZ$ 是一个\omterm{def:b3:clt:gaussianvector}{高斯向量}
        （见\cref{def:b3:clt:gaussianvector}），均值为 $0$，
        协方差为 $\Sigma$，且 $\gamma' = \E\norm N^3 <
        \infty$；并证明 $G_n = \frac{N_1 + \dots +
        N_n}{\sqrt n}$ 服从 $\mathcal N(0, \Sigma)$，
        对\omterm{def:b3:probability:independence}{独立}同分布的副本 $N_i$，\emph{恰好}成立。
  \item （二元泰勒展开）设 $f \colon \R^2 \to
        \R$ 为 $\mathcal C^3$ \omterm{def:b3:representations:character}{类函数}，且 $M_3 =
        \max_{\abs\alpha = 3}\sup\abs{\partial^\alpha f} <
        \infty$，证明
        \[
        \Bigl|f(w + h) - f(w) - \langle\nabla f(w),
        h\rangle - \tfrac12\langle h, D^2f(w)\,h\rangle
        \Bigr| \leq \frac{M_3}6\,\bigl(\abs{h_1} +
        \abs{h_2}\bigr)^3 \leq \frac{\sqrt2\,M_3}3\,
        \norm h^3
        \]
        \emph{（考察 $t \mapsto f(w + th)$ 在
        $\intcc01$ 上的情形）}。
  \item （$\R^2$ 中的中心极限定理）用替换方案处理
        向量混合量 $H_i$：证明一阶与二阶项
        相消（均值与协方差相配），裂项求和，
        并仿照第 5 题加强
        （紧性由 $\E\norm{T_n}^2 =
        \operatorname{tr}\Sigma$ 得出；光滑化现在
        在 $\R^2$ 中进行，见\cref{thm:b3:lp:regularization}）
        推出结论：对每个有界连续函数 $f \colon
        \R^2 \to \R$，
        \[
        \E\,f\Bigl(\frac{X_1 + \dots + X_n}{\sqrt n}\Bigr)
        \longrightarrow \E\,f(N), \qquad N \sim \mathcal
        N(0, \Sigma) :
        \]
        这即 \cref{thm:b3:clt:multiclt} 在 $2$ 维情形下，
        对光滑的 $f$ 给出速率，且无需傅里叶分析。
  \item （克拉默--沃尔德，以及联合涨落）由此推出
        $\langle t, \frac{S_n}{\sqrt n}\rangle
        \Rightarrow \mathcal N(0, t^{\mathsf T}\Sigma t)$
        对每个固定的 $t \in \R^2$ 成立。应用：对
        \omterm{def:b3:probability:independence}{独立}同分布的实\omterm{def:b3:probability:space}{随机变量} $(\xi_i)$，中心化，$\E\xi_1^2 = 1$，
        $\E\xi_1^6 < \infty$（使得第五部分可用于
        $V_i = (\xi_i, \xi_i^2 - 1)$），证明
        \[
        \frac1{\sqrt n}\Bigl(\sum_{i\leq n}\xi_i,\
        \sum_{i\leq n}(\xi_i^2 - 1)\Bigr) \Longrightarrow
        \mathcal N\Bigl(0, \begin{pmatrix} 1 & \E\xi_1^3\\
        \E\xi_1^3 & \E\xi_1^4 - 1\end{pmatrix}\Bigr) :
        \]
        经验均值与经验二阶矩
        联合呈高斯涨落——在极限下相互\omterm{def:b3:probability:independence}{独立}，当且仅当
        $\E\xi_1^3 = 0$ 时成立
        （见\cref{thm:b3:clt:gaussianvector}）。
\end{enumerate}

\medskip
\noindent\textbf{第六部分——德尔塔方法。}
\begin{enumerate}[resume]
  \item 设 $(\hat\theta_n)$ 为\omterm{def:b3:probability:space}{随机变量}，满足
        $\sqrt n(\hat\theta_n - \theta) \Rightarrow
        \mathcal N(0, \sigma^2)$，其中 $\theta$ 为实参数，
        设 $g$ 在
        $\theta$ 处可导。证明\emph{德尔塔方法}：
        \[
        \sqrt n\bigl(g(\hat\theta_n) - g(\theta)\bigr)
        \Longrightarrow \mathcal N\bigl(0,
        g'(\theta)^2\sigma^2\bigr)
        \]
        \emph{（写 $g(x) - g(\theta) = (g'(\theta) +
        \eta(x))(x - \theta)$，其中 $\eta \to 0$ 在
        $\theta$ 处；先证 $\hat\theta_n \to \theta$，
        依概率有 $\eta(\hat\theta_n) \to 0$，最后
        用斯卢茨基定理、\cref{exo:b3:clt:8} 和
        \cref{exo:b3:clt:4}(b)）}。
  \item 应用。(a) 设 $(\xi_i)$ 为\omterm{def:b3:probability:independence}{独立}同分布的实值序列，
        其均值为 $\mu$、方差为 $\sigma^2$，并记 $\bar X_n =
        \frac1n\sum_{i\leq n}\xi_i$：证明 $\sqrt n(\bar
        X_n^2 - \mu^2) \Rightarrow \mathcal N(0,
        4\mu^2\sigma^2)$ 当 $\mu \neq 0$ 时成立，并说明
        $\mu = 0$ 时，正确的结论需在另一尺度上给出：
        $n\bar X_n^2 \Rightarrow \sigma^2N^2$，其中
        $N \sim \mathcal N(0,1)$（求出极限的
        分布函数）。(b)（方差
        稳定化）。设 $\hat p_n$ 为
        $\mathcal B(1, p)$ 样本的成功频率，$p \in \intoo01$：
        证明 $g(p) = \arcsin\sqrt p$ 满足
        \[
        \sqrt n\,\bigl(g(\hat p_n) - g(p)\bigr)
        \Longrightarrow \mathcal N\Bigl(0, \frac14\Bigr)
        \]
        \emph{无论} $p$ 取何值——渐近误差
        范围不依赖未知参数；可与
        \cref{ex:b3:clt:confidence} 比较。
\end{enumerate}

\medskip
\noindent\textbf{第七部分——泊松分布，用同样的方法：
勒卡姆定理。}替换法还揭示了第二类普适类：
许多\omterm{def:b3:probability:independence}{独立}的\emph{稀有}事件之和。对
$\N$ 上的分布，正确的距离是\emph{全变差}，
\[
d_{\mathrm{TV}}(\mu, \nu) = \sup_{A\subseteq\N}\,
\abs{\mu(A) - \nu(A)} .
\]
\begin{enumerate}[resume]
  \item 证明 $d_{\mathrm{TV}}(\mu, \nu) =
        \frac12\sum_{k\geq0}\abs{\mu(\{k\}) -
        \nu(\{k\})}$，并证明耦合界：对
        \emph{任意}\omterm{def:b3:probability:space}{随机变量}对 $(X, Y)$，其
        分布分别为 $\mu$ 与 $\nu$，且定义在同一空间上，
        $d_{\mathrm{TV}}(\mu, \nu) \leq \P(X \neq Y)$。
  \item 对 $p \in \intoo01$，精确计算：
        \[
        d_{\mathrm{TV}}\bigl(\mathcal B(1, p), \mathcal
        P(p)\bigr) = p\bigl(1 - \eu^{-p}\bigr) \leq p^2 .
        \]
  \item （勒卡姆，交换法）设 $X_i \sim \mathcal B(1,
        p_i)$ 与 $Y_i \sim \mathcal P(p_i)$，共 $2n$
        个变量相互\omterm{def:b3:probability:independence}{独立}；$S = X_1 + \dots + X_n$，并
        回忆 $Y_1 + \dots + Y_n \sim \mathcal P(\lambda)$，
        其中 $\lambda = \sum_ip_i$（见\cref{exo:b3:clt:1}）。
        每次交换一个坐标，构造整数型
        混合量 $H_i = Y_1 + \dots + Y_i + X_{i+1} + \dots
        + X_n$：对任意 $A \subseteq \N$，证明
        \[
        \abs{\P(H_{i-1} \in A) - \P(H_i \in A)} \leq
        d_{\mathrm{TV}}\bigl(\mathcal B(1, p_i), \mathcal
        P(p_i)\bigr),
        \]
        并得出\emph{勒卡姆不等式}：
        \[
        d_{\mathrm{TV}}\bigl(\text{ 的分布 } S,\ \mathcal
        P(\lambda)\bigr) \leq \sum_{i=1}^np_i^2 .
        \]
  \item 收获。(a) 对 $p_i = \frac\lambda n$，
        误差界为 $\frac{\lambda^2}n$——稀有事件律
        （见\cref{exo:b3:clt:5}）升级为显式速率，
        对所有事件一致，且
        对不等的 $p_i$ 同样有效。(b) $500$ 封信被投递，
        各封信\omterm{def:b3:probability:independence}{独立}地误投，概率为
        $\frac1{500}$：给出
        参数为 $1$ 的泊松模型的误差界，并估计
        一封信都不误投的概率。(c) 结束本题：比较
        此处讨论的两类普适类——
        高斯类（许多微小且分散的贡献；
        两阶矩匹配；泰勒展开）与
        泊松类（许多稀有贡献；一阶矩
        匹配；精确的全变差耦合）——以及
        二者背后统一的替换法。

  \item （相对误差与对数变换）设 $(X_n)$
        为\omterm{def:b3:probability:independence}{独立}同分布的正值序列，均值 $\mu > 0$，方差
        $\sigma^2$，并记 $\bar X_n$ 为经验均值。
        用德尔塔方法证明下式：
        \[
        \sqrt n\,\bigl(\ln\bar X_n - \ln\mu\bigr)
        \Longrightarrow
        \mathcal N\Bigl(0,\ \frac{\sigma^2}{\mu^2}\Bigr) :
        \]
        $\ln\bar X_n$ 的渐近参数是
        \emph{变异系数} $\sigma/\mu$——
        相对且与尺度无关的误差。由此导出 $95\%$
        $\mu$ 的乘性置信区间，
        形式为 $\bar X_n\cdot\eu^{\pm1.96\,\sigma/(\mu\sqrt
        n)}$，并说明何时它优于
        加性区间。
  \item （三阶矩决定误差）设 $X \sim$
        中心化的伯努利（$p$）分布，计算 $\E\bigl[(X -
        p)^3\bigr] = p(1-p)(1-2p)$。利用第四部分的分析
        （交换误差由三阶矩主导），解释
        为何 $\mathcal B(n, p)$ 的正态近似
        在 $p \neq \frac12$ 时不对称：一侧过估，
        另一侧欠估；并解释为何
        $p = \frac12$ 时匹配矩的收敛速度更快。
        数值验证偏度的符号：在 $\mathcal
        B(20, 0.1)$ 与 $\mathcal N(2, 1.8)$ 的情形下，
        比较 $\P(S = 0) = 0.9^{20}$ 与\omterm{ex:b3:probability:laws}{高斯分布}在
        $\intoo{-\infty}{0.5}$ 上的质量。
\end{enumerate}
\end{problem}
